Esta tese apresenta o HySail como uma proposta de armazenamento distribuído com foco em recuperabilidade, integridade e auditoria algébrica. O texto organiza a discussão em torno da introdução do protocolo, de seus fundamentos matemáticos, de seu funcionamento operacional, da integração com blockchain e dos resultados observados na implementação atual. São discutidos os papéis dos Online Codes, da representação de blocos em corpo de Galois, da divisão polinomial em $GF(2)$ e da função \texttt{gf2\_poly\_mod} na construção dos autenticadores locais. Também é analisada a utilização de contratos inteligentes em Ethereum como camada de coordenação e rastreabilidade para registro de arquivos, provedores e solicitações de reconstrução. Os resultados indicam que a base conceitual do Hy-SAIL permanece presente no repositório, embora a codificação corrente se aproxime mais de um esquema inspirado em LT Code do que de uma implementação integral de Online Code.
